\documentclass[11pt,a4paper]{article}
\usepackage[utf8]{inputenc}
\usepackage[french]{babel}
\usepackage[T1]{fontenc}
\usepackage{lmodern}
\usepackage[margin=1.5cm]{geometry}

\begin{document}

% En-tête
\begin{center}
    {\Huge \textbf{Samy Martin}} \\[0.2cm]
    {\Large \textbf{Programmeur Moteur 3D}} \\[0.1cm]
    samy.martin@email.com \ | \ 06 67 10 57 72
\end{center}

\vspace{0.3cm}

% Résumé / Profil
\section*{Profil Professionnel}
\noindent
Ingénieur spécialisé dans le développement de jeux vidéo et la programmation 3D. Passionné par l'optimisation des performances en temps réel et l'architecture de moteurs de jeu, je possède une excellente maîtrise du C++ et des outils de l'industrie (Unreal Engine). Mon approche mathématique et ma rigueur me permettent de repousser les limites visuelles et techniques des plateformes modernes.

\vspace{0.3cm}

% Compétences
\section*{Compétences Techniques}
\noindent
\textbf{Optimisation GPU} $\bullet$ \textbf{C++} $\bullet$ \textbf{Vulkan} $\bullet$ \textbf{Git} $\bullet$ \textbf{OpenGL} $\bullet$ \textbf{Mathématiques 3D} $\bullet$ \textbf{Perforce} $\bullet$ \textbf{Unreal Engine 5}

\vspace{0.3cm}

% Expériences
\section*{Expériences Professionnelles}
\noindent \textbf{Programmeur Moteur 3D / Confirmé} --- \textit{Asobo Studio} \hfill \textbf{09/2021 à Aujourd'hui} \\
Optimisation de la gestion de la mémoire et réduction du temps de chargement des niveaux de 40\% grâce au streaming asynchrone des assets 3D. Portage complet du moteur de jeu vers les consoles nouvelle génération (PS5, Xbox Series X), en exploitant au maximum les architectures SSD. Mise en place d'une architecture réseau fiable pour le mode multijoueur (client-serveur avec prédiction côté client et compensation du lag). Intégration d'un moteur physique multi-threadé permettant de simuler des destructions de décors complexes à grande échelle sans perte de framerate. \\[0.3cm]
\noindent \textbf{Programmeur Moteur 3D} --- \textit{Dontnod Entertainment} \hfill \textbf{02/2019 à 08/2021} \\
Implémentation de shaders personnalisés et de systèmes de particules avancés avec Niagara pour améliorer la qualité visuelle globale du jeu. Développement d'outils internes pour l'éditeur de niveau, augmentant de 30\% la productivité de l'équipe de Level Design au quotidien. Développement de mécaniques de gameplay en C++ et optimisation du rendu graphique sur Unreal Engine pour un titre AAA, garantissant 60 FPS constants. Intégration d'outils de télémétrie in-game pour récolter les données de parcours des joueurs et ajuster la difficulté dynamiquement. \\[0.3cm]
\noindent \textbf{Alternance — Assistant Programmeur Moteur 3D} --- \textit{Arkane Studios} \hfill \textbf{12/2018 à 01/2019} \\
Profilage CPU/GPU intensif à l'aide de PIX et RenderDoc pour identifier et corriger les goulots d'étranglement avant la phase de certification console. Développement d'un système de cinématiques temps réel, permettant de synchroniser les animations faciales avec le doublage audio. Création d'un système d'intelligence artificielle (Behavior Trees) pour les PNJ, rendant leurs réactions plus réalistes et adaptatives face au joueur. \\[0.3cm]


\vspace{0.3cm}

% Formations
\section*{Formation \& Diplômes}
\noindent \textbf{Diplôme d'Ingénieur en Programmation 3D} \hfill \textbf{2016 - 2019} \\
\textit{Cnam ENJMIN} \\[0.3cm]
\noindent \textbf{DUT en Programmation 3D} \hfill \textbf{2014 - 2016} \\
\textit{IUT Informatique et Imagerie} \\[0.3cm]


\end{document}